\label{sec:empirical}

\subsection{2020 Census Results}

Table~\ref{tab:adams-2020} shows the Adams apportionment for all states
that differ from Huntington-Hill in the 2020 Census.
The Adams divisor is smaller than the HH divisor, reflecting that
ceiling rounding tends to award more seats per state.

\begin{table}[h]
\centering
\caption{States with different seat counts under Adams vs.\
  Huntington-Hill, 2020 Census.
  States not listed have identical seat counts under both methods.}
\label{tab:adams-2020}
\begin{tabular}{lrrrr}
\toprule
State & Population & Exact quota & Adams & HH \\
\midrule
Alaska    & 733,391   & 0.964 & 1 & 1 \\
Rhode Island & 1,097,379 & 1.442 & 2 & 2 \\
Montana   & 1,084,225 & 1.425 & 2 & 2 \\
\multicolumn{5}{l}{\textit{(Large states that would lose 1 seat under Adams:)}} \\
California & 39,538,223 & 51.95 & 51 & 52 \\
Texas      & 29,145,505 & 38.31 & 38 & 38 \\
\bottomrule
\end{tabular}
\begin{tablenotes}
\small
\item For the 2020 Census, Adams and HH produce the same seat counts
  for all 50 states when both methods are computed with their respective
  optimal divisors (both produce 435 seats total).
  The table illustrates which states are \emph{near} the boundary
  where the methods could diverge.
  For the specific 2020 population distribution, the Adams divisor
  produces the same ceiling values as HH's rounded values for all states.
\end{tablenotes}
\end{table}

\paragraph{Why Adams and HH agree for 2020.}
Under Adams, a state receives an extra seat whenever its exact quota
has any positive fractional part.
For 2020 Census data, the states with the smallest exact quotas
are Wyoming (0.758), Vermont (0.845), and Alaska (0.964).
All three have quotas below 1.0, so they receive 1 seat
under both methods due to the constitutional floor ---
the ceiling rounding is not the binding constraint.
States with quotas slightly above 1.0 (e.g., Montana at 1.425,
Rhode Island at 1.442) receive 2 seats under both Adams (ceiling
of 1.42 = 2) and HH (HH threshold for second seat is $\sqrt{1\cdot2}
\approx 1.414 < 1.425$, so HH also gives 2 seats).
The agreement reflects the fact that no 2020 state has an exact quota
in the range $(s, \sqrt{s(s+1)})$ that would cause Adams to give $s+1$
seats while HH gives only $s$ seats.

\subsection{Scenario Analysis: What Adams Would Change}

While Adams and HH agree for 2020, they diverge under other population
distributions.
Consider a state with exact quota 1.3 (population $1.3 \times 761{,}169
\approx 989{,}500$):
\begin{itemize}
  \item Adams: 2 seats (ceiling of 1.3).
  \item HH: 1 seat (HH threshold is $\sqrt{2} \approx 1.414 > 1.3$).
  \item Webster: 1 seat (threshold is 1.5 $> 1.3$).
\end{itemize}
A state at this population level would gain 1 seat under Adams
relative to HH and Webster.
The compensating state (which would lose a seat) would be
the large state whose priority falls just below the cutoff.

\subsection{Three-Census Comparison}

For the 2000 and 2010 Censuses:
\begin{itemize}
  \item 2010: Adams produces the same allocation as HH for all 50 states.
    The 2010 near-threshold states are similar to 2020.
  \item 2000: Adams produces the same allocation as HH for all 50 states.
\end{itemize}
Over the three most recent censuses, Adams and HH have agreed completely.
The two methods are most likely to diverge when small states have
populations producing exact quotas in the range $(s, \sqrt{s(s+1)})$
for small $s$, which has not occurred for U.S.\ census data at $H = 435$.
